\documentclass[10pt,a4paper]{amsart}
\usepackage[margin=19mm]{geometry}
\usepackage{amsmath,amssymb,mathtools}
\usepackage[scr=rsfs,cal=euler]{mathalfa}
\usepackage{xfrac}
\usepackage[unicode,hidelinks,bookmarks=false]{hyperref}
\newcommand{\ie}{i.e.\ }
\newcommand{\cf}{cf.\ }
\newcommand{\resp}{respectively\ }
\newcommand{\Subsection}{\subsection}
\providecommand{\nobreakdash}{\nobreak\hbox{-}}
\setlength{\emergencystretch}{2em}
\multlinegap0pt
\usepackage{mathtools}
\usepackage{amssymb}
\usepackage{booktabs}
\usepackage[shortlabels]{enumitem}
\usepackage{tikz}
\usepackage{dsfont}
\usepackage{multirow}
\newtheorem{lemma}{Lemma}
\DeclareMathOperator{\Lk}{Lk}
\DeclareMathOperator{\Pic}{Pic}
\title{The characterization of $(n-1)$-spheres with $n+4$ vertices having maximal Buchstaber number}
\author{Suyoung Choi, Hyeontae Jang, Mathieu Vall\'ee}
\date{Source: arXiv:2301.00806v3 (CC BY 4.0)}
\begin{document}
\maketitle

\noindent\emph{Source and licence.} The characterization of $(n-1)$-spheres with $n+4$ vertices having maximal Buchstaber number. Version arXiv:2301.00806v3 (CC BY 4.0), distributed by arXiv under the Creative Commons Attribution 4.0 International licence, http://creativecommons.org/licenses/by/4.0/, as recorded in the arXiv licence metadata for this exact version and shown on the abstract page https://arxiv.org/abs/2301.00806v3. Full licence notice: \texttt{licenses/arxiv-2301.00806-CC-BY-4.0.txt}.\par\smallskip

\noindent\textit{Editorial scope.} Counted unit: the article's lemma lemma (source0.tex lines 970--974). The article's complete original proof of it is delivered with it (source0.tex lines 975--986, one \texttt{proof} environment of 12 lines). No mathematics is written, repaired or re-derived: the statement and the proof are the article's own lines, in the article's own notation, and no retained line is substituted or re-flowed.  No further source-local context is needed: every cross-reference of every retained line resolves inside the two delivered spans.  Nothing was omitted from any retained span, so the package carries no omission record.  No retained line contains a citation, so the wrapper carries no bibliography and no bibliography entry.  No retained line contains a figure environment, an image-inclusion command or drawing code, so no image is delivered and the package declares no asset dependency.  Editorial formatting only, in addition: an AMS \texttt{amsart} preamble that restates the article's own declarations and macros used by the retained lines, namely the environments \texttt{lemma} on separate counters, together with the standard packages those lines need.  The retained environments print in this document as Lemma 1 in order of appearance, whereas the article prints the same items with the numbering of its own full document.  The complete native source of the article as distributed in the arXiv e-print for this exact version is bundled unchanged as \texttt{sources/135631/source0.tex}.
\begin{lemma}
	Let $K$ be a seed of Picard number~$4$.
	Assume that $K$ has a vertex $v$ such that $\Pic(\Lk_K(v)) = 3$ and there exist two vertices $v_1$ and $v_2$ of $\Lk_K(v)$ such that every facets of $\Lk_K(v)$ contains either $v_1$ or $v_2$.
	Then, there is a vertex of $K$ whose link has Picard number~$4$.
\end{lemma}

\begin{proof}
	Let $\{v_1\} \cup \sigma$ be a facet without $v_2$.
	There is one more facet containing $\sigma$ since it is a ridge of the PL~sphere $\Lk_K(v)$.
	By the assumption, it must be $\{v_2\} \cup \sigma$.
	That shows every vertex in $\Lk_K(v)$ forms an edge with both $v_1$ and $v_2$.
	
	Let $w$ be the vertex not in $\Lk_K(v)$.
	If $v_1 \in \Lk_K(w)$, then $\Lk_K(v_1)$ has Picard number~$4$.
	If both $v_1, v_2 \not\in \Lk_K(w)$, then $\Lk_K(w)$ is an $(n-2)$-dimensional PL~sphere with $n$ vertices.
    This means that $\Lk_K(w) = \partial\Delta_{n-1}$.
    Then if $w^\prime$ is a vertex of $\Lk_K(w)$ other than $w$, then $\Pic(\Lk_K(w^\prime)) = 4$.
\end{proof}

\end{document}
